% ARCHIVED DRAFT: independently audited by root and incorporated on 2026-11-19 16:28 UTC.
% The assembled paper is authoritative; this file retains drafting provenance.
\begin{saresult}{B.069}{Proved under stated hypotheses}{Power-generated maximal ideals and surjective Frobenius}
Let $R$ be a nonzero commutative ring. Suppose that, at each maximal ideal
$\mathfrak m$ with finite residue field, the local maximal ideal
$\mathfrak n=\mathfrak mR_{\mathfrak m}$ satisfies
\[
             \mathfrak n=(a^N:a\in\mathfrak n)
             \quad\text{for every integer }N\ge1.          \tag{B69.1}
\]
Then for every unital inclusion $R\subset S$ the following conditions
are equivalent:
\[
\begin{gathered}
 \text{the polynomial co-operations exist uniquely};\\
 \theta_2\text{ is injective};\qquad C=R;\\
 (R:_S D_d(R))=R\quad\text{for every }d\ge1;\\
 \operatorname{Ann}_{S/R}(D_d(R))=0\quad\text{for every }d\ge1;\\
 \Omega(R,S)=R[X],                                        \tag{B69.2}
\end{gathered}
\]
where $C$ is the actual coefficient ring of B.064. Every positive case
has all multivariable polynomial tensor isomorphisms.

In particular this class includes every commutative $\mathbb F_p$-algebra
with surjective Frobenius, with arbitrary nilpotents and without
Noetherian assumptions. Surjectivity, rather than injectivity, of
Frobenius is what this assertion requires.
\end{saresult}
\begin{proof}
We first compute the determinant ideals locally, and do not assume
localization of integer-valued polynomial modules. For any multiplicative
set $U$, one has
\[
                   D_d(U^{-1}R)=U^{-1}D_d(R).              \tag{B69.3}
\]
The forward inclusion follows by giving a finite tuple of nodes a common
denominator $u\in U$: its Vandermonde determinant is the determinant of
nodes in $R$, divided by $u^{d(d+1)/2}$. The reverse inclusion is immediate.

Let $(A,\mathfrak n)$ be a local ring satisfying (B69.1), with residue
field of cardinality $q$. When $d<q$, lifts of $d+1$ distinct residue
elements have unit Vandermonde determinant, so $D_d(A)=A$. When $d\ge q$,
any $d+1$ residues repeat, so $D_d(A)\subset\mathfrak n$.
For the reverse inclusion use the polynomial $P_d$ of B.064. At any
$a\in\mathfrak n$ its explicit factors give
\[
 P_d(a)=a^{e_d}u_a,\qquad
 e_d=\frac{d(d-1)(d+1)}6,\qquad u_a\in A^\times.            \tag{B69.4}
\]
Indeed every difference between powers is a power of $a$ times a factor
$a^j-1$, which is a unit. The exponent is the sum of the exponents from
the differences with zero and the smaller exponents from pairs of
positive powers. Thus $a^{e_d}\in D_d(A)$ for every $a\in\mathfrak n$.
Equation (B69.1), with $N=e_d\ge1$, proves
\[
 D_d(A)=
 \begin{cases}A,&d<q,\\ \mathfrak n,&d\ge q.\end{cases}     \tag{B69.5}
\]
Also (B69.1) with $N=2$ gives $\mathfrak n^2=\mathfrak n$.
At an infinite-residue local ring all $D_d(A)$ equal $A$, by distinct
residue nodes, without any hypothesis on its maximal ideal.

Suppose now that $\theta_2$ is injective. B.064 makes $R\to C$ an
injective ring epimorphism. Localizing this inclusion and the epimorphism
at any maximal ideal $\mathfrak m$ gives $A=R_{\mathfrak m}\subset
B=C_{\mathfrak m}$. For each $b\in B$, some $D_d(A)$ multiplies $b$ into
$A$, because $b$ comes from one $C_d$ and (B69.3) applies. If the residue
field is infinite this already says $B=A$.

In the finite-residue case, (B69.5) says
$\mathfrak nB\subset A$. As an $A$-ideal it contains $\mathfrak n$, so it
equals $\mathfrak n$ or $A$. The latter is impossible. Otherwise write
$1=\sum_{i=1}^s a_i b_i$ with $a_i\in\mathfrak n$, $b_i\in B$. For every
$a\in\mathfrak n$ the coefficients $ab_i$ lie in $A$, so
$a=\sum_i a_i(ab_i)$ proves that $\mathfrak n$ is finitely generated.
But $\mathfrak n^2=\mathfrak n$ and Nakayama's elementary finite-generator
lemma then give $\mathfrak n=0$, contradicting $1\in\mathfrak nB$.
Hence
\[
                         \mathfrak nB=\mathfrak n.         \tag{B69.6}
\]
Base change of the epimorphism by $A\to A/\mathfrak n$ is therefore an
injective epimorphism from the field $k=A/\mathfrak n$ into
$B/\mathfrak n$. Such an inclusion is an isomorphism: if $x$ in a
$k$-algebra is linearly independent of $1$, then
$x\otimes1-1\otimes x$ is nonzero over $k$, contradicting the epimorphism
condition. Thus $B/\mathfrak n=k$ and $B=A+\mathfrak n=A$.

We have proved $C_{\mathfrak m}=R_{\mathfrak m}$ at every maximal ideal.
The module $C/R$ therefore vanishes, with no finite-generation hypothesis:
a nonzero element has a proper annihilator contained in a maximal ideal,
and survives localization there. Hence $C=R$. B.064 proves all remaining
equivalences in (B69.2) and the positive tensor isomorphisms.

Finally suppose $R$ is an $\mathbb F_p$-algebra with surjective Frobenius.
This property survives localization. For $a$ in a local maximal ideal
and any $N$, choose $e$ with $p^e\ge N$ and a $p^e$-th root $b$ of $a$.
Then $b$ also lies in the maximal ideal, since it is prime, and
$a=b^{p^e}=b^N b^{p^e-N}$. Every element of that maximal ideal lies in the
ideal generated by $N$-th powers of its elements, proving (B69.1).
The reverse containment is automatic. Nilpotents cause no change in
this argument.
\end{proof}
